\chapter{Empirical Evaluation}\label{ch:evaluation}

The evaluation is designed around claims that can be supported by the implementation.
Motivo is not evaluated as a general measure of musical creativity, and the thesis does not include a user study.
Instead, this chapter evaluates engineering properties: whether the compiler can expand compact source programs into
larger MIDI event streams, how quickly it compiles a small benchmark corpus on one local machine, what safety limits the
current implementation imposes, and how the software project is tested.

\section{Methodology}\label{sec:evaluation-methodology}

The benchmark was run locally with a Release build of \texttt{motivoc}.
The binary was built outside the repository at \texttt{/tmp/motivo-thesis-bench-build} using CMake 4.3.2 and Apple
clang 21.0.0.
Measurements were taken on an Apple M3 Pro machine with 18 GB of RAM\@.
The timing tool was \texttt{hyperfine} 1.20.0 with shell execution disabled, 5 warmup runs, and 50 measured
runs~\cite{Hyperfine_2026}.
Each measured command compiled one \texttt{.motivo} source file and wrote the resulting MIDI file, so the
results include
process startup, parsing, semantic analysis, lowering, MIDI serialization, and file output.

The benchmark corpus was synthetic.
It was built to exercise source-to-event expansion rather than to produce musically refined pieces.
This distinction matters.
A loop that generates 20,000 notes is useful for measuring expansion and compilation cost, but it does not prove that
the generated score is artistically interesting.
The benchmark therefore supports the engineering claim that Motivo can generate many symbolic events from
compact source rules within the tested limit.

A defining characteristic of the Motivo compiler is its ``zero-runtime'' architecture, which shifts the
computational burden entirely to the compilation phase.
To evaluate the viability of this approach, it is essential to analyze the compiler's execution speed, particularly as
the complexity of the musical composition scales.

\subsection{Theoretical Complexity}\label{detail-subsec:theoretical-complexity}

The compilation pipeline is composed of several sequential passes, each with a distinct computational complexity:

\noindent \textbf{Frontend and Semantic Analysis.} The complexity of parsing and type-checking is proportional to the
number of lines of source code, $O(S)$, where $S$ is the size of the abstract syntax tree.
Because these passes do not unroll loops or instantiate patterns multiple times, they remain extremely fast
regardless of the composition's length.

\noindent \textbf{Lowering Engine.} The unrolling process has a time complexity of $O(E)$, where $E$ is the total number
of generated \texttt{NoteEvent}s. A simple \texttt{loop (1000)} statement will execute its body 1,000 times during
this phase, meaning the lowering time is proportional to the density of the final timeline, not the brevity of the
source code.

\noindent \textbf{MIDI Backend.} The backend must sort the absolute events by their tick values before computing
delta-times.
This stable sort introduces a time complexity of $O(E \log E)$.
Serializing the events into the binary file requires a final linear pass, $O(E)$.

Therefore, the overall time complexity of the compilation process is bounded by the sorting phase in the backend:
$O(E \log E)$.

\section{Compilation Time}\label{sec:compile-time}

Table~\ref{tab:benchmark-results} shows the measured results.
The fastest case and the realistic example are close to 2 ms.
The largest synthetic cases are around 6 ms.
Because these times are small, they should be interpreted as local measurements on the stated machine rather than as a
portable performance result.
They are still useful for the thesis because they show that the current native compiler is fast enough to support an
interactive compile loop in this local setup.

\begin{table}[htbp]
  \centering
  \footnotesize
  \renewcommand{\arraystretch}{1.25}
  \begin{tabular}{lrrrrr}
    \toprule
    \textbf{Program} & \textbf{LOC} & \textbf{Note events} & \textbf{MIDI size} & \textbf{Mean ms} & \textbf{SD ms} \\
    \midrule
    \texttt{single\_loop\_100} & 7 & 100 & 856 B & 2.02 & 0.31 \\
    \texttt{single\_loop\_1000} & 7 & 1,000 & 8,056 B & 2.14 & 0.31 \\
    \texttt{single\_loop\_10000} & 7 & 10,000 & 80,056 B & 3.96 & 0.24 \\
    \texttt{single\_loop\_20000} & 7 & 20,000 & 160,056 B & 6.03 & 0.15 \\
    \texttt{chord\_loop\_5000} & 7 & 20,000 & 160,056 B & 5.90 & 0.17 \\
    \texttt{drum\_groove\_500} & 18 & 10,000 & 88,053 B & 6.39 & 0.19 \\
    \texttt{example\_copy} & 77 & 188 & 1,751 B & 2.04 & 0.18 \\
    \bottomrule
  \end{tabular}
  \caption{Local Release-build compilation benchmark.
  Times include MIDI file writing and are based on 50 measured runs after 5 warmups.}\label{tab:benchmark-results}
\end{table}

\begin{figure}[htbp]
  \centering
  \begin{tikzpicture}
    \begin{axis}[
        ybar,
        width=0.95\textwidth,
        height=7cm,
        ymin=0,
        ylabel={Mean compile time (ms)},
        symbolic x coords={100,1000,10000,20000,chord,drums,example},
        xtick=data,
        x tick label style={rotate=35, anchor=east, font=\scriptsize},
        nodes near coords,
        nodes near coords style={font=\scriptsize},
        bar width=12pt,
        enlarge x limits=0.08,
        ymajorgrids=true,
        grid style={draw=motivoInk!12}
      ]
      \addplot[fill=motivoBlue!70, draw=motivoBlue!90] coordinates {
        (100,2.02)
        (1000,2.14)
        (10000,3.96)
        (20000,6.03)
        (chord,5.90)
        (drums,6.39)
        (example,2.04)
      };
    \end{axis}
  \end{tikzpicture}
  \caption{Mean compile times for the local benchmark corpus.}\label{fig:benchmark-time}
\end{figure}

The near-linear growth in the single-loop cases is consistent with the implementation.
Parsing and semantic analysis see a small source program, while lowering and MIDI writing scale with the number of
generated events.
The backend also sorts events by tick before computing delta-times.
The result is still small at the current safety limit, but the implementation should not be described as unbounded.

\subsection{Execution Speed and Real-Time Viability}\label{detail-subsec:execution-speed-and-real-time-viability}

Because the compiler is implemented as a native C++23 program without reliance on garbage collection or dynamic dispatch
for AST traversal, the constant factors in the time complexity are small in the current measurements.
In practice, the compiler generates compositions containing event streams up to the current implementation
limit in a fraction of a second.

These local measurements are relevant to Motivo Studio because the editor workflow depends on compilation
remaining quick enough for interactive use.
The interactive playground environment (discussed in previous chapters) relies on the compiler's ability to re-compile
the entire source file during interactive compile actions.
Because the execution time for typical musical pieces (ranging from hundreds to thousands of events) sits
well below the latency usually noticed during ordinary editor interaction, the composer experiences
interactive auditory and visual feedback without using an interpreted runtime during playback.

\section{Source-to-Event Expansion}\label{sec:compression}

A major objective of Motivo is to provide a highly expressive interface for music composition,
reducing the manual effort required to construct complex arrangements.
We evaluate this expressiveness by analyzing the ``compression ratio''---the difference between the amount of
code written
by the composer and the number of low-level MIDI events generated by the compiler.

\subsection{Algorithmic Conciseness}\label{detail-subsec:algorithmic-conciseness}

Traditional MIDI sequencers and digital audio workstations require the composer to manually place every note on a grid.
In contrast, Motivo allows composers to define structural rules and relationships.

Consider the following drum pattern from the \texttt{example.motivo} file:
\begin{lstlisting}[language=Motivo, caption={A generative drum pattern using control flow.},label={lst:evaluation-drum-pattern}]
track percussion using drums {
    pattern groove() {
        for (let i = 0; i < 16; i = i + 1) {
            play hihat from i;
            if (i % 8 == 0) { play kick from i; }
            if (i % 8 == 4) { play snare from i; }
        }
    }
    loop (3) { play groove(); }
}
\end{lstlisting}

This 10-line block of code encapsulates a complex rhythm.
The \texttt{for} loop defines the underlying hi-hat grid, while the modulo-based \texttt{if} conditions procedurally
generate the kick and snare syncopation.
The \texttt{loop (3)} statement then instantiates this 16-beat pattern three times across the timeline.

\subsection{The Compression Ratio}\label{detail-subsec:the-compression-ratio}

To achieve the same result in a traditional MIDI editor, the composer would have to manually draw 48 hi-hat notes, 6
kick drum notes, and 6 snare drum notes---a total of 60 individual note events, which translates to 120 low-level MIDI
messages (\texttt{Note-On} and \texttt{Note-Off}).

Motivo abstracts these 120 MIDI messages into just a handful of declarative rules.
The resulting ``compression ratio'' is exceedingly high.
As the composition grows---for example, looping the groove 32 times instead of 3---the size of Motivo source
code remains
completely unchanged, while the output scales to hundreds or thousands of events.

This expressiveness validates the core thesis: by shifting from a linear, event-based representation to a structural,
programming-based representation, the manual overhead of repetitive event placement can be reduced.
The composer can focus on higher-level musical structure while the lowering engine handles repetitive
temporal event placement.

The main measured expressiveness result is source-to-event expansion.
The 20,000-event single-loop program uses 7 lines of source.
The chord-loop program also uses 7 lines and generates 20,000 note events because each loop iteration emits a
four-note chord.
The 18-line drum groove generates 10,000 note events by combining a 16-step bar pattern with a 500-iteration loop.
These examples demonstrate the intended effect of the language: repeated symbolic structure remains compact in source
form while the compiler expands it into explicit event data.

\begin{figure}[htbp]
  \centering
  \begin{tikzpicture}
    \begin{axis}[
        ybar,
        width=0.95\textwidth,
        height=7cm,
        ymin=0,
        ylabel={Note events generated per source line},
        symbolic x coords={100,1000,10000,20000,chord,drums,example},
        xtick=data,
        x tick label style={rotate=35, anchor=east, font=\scriptsize},
        nodes near coords,
        nodes near coords style={font=\scriptsize, rotate=90, anchor=west},
        bar width=12pt,
        enlarge x limits=0.08,
        ymajorgrids=true,
        grid style={draw=motivoInk!12}
      ]
      \addplot[fill=motivoGreen!70, draw=motivoGreen!90] coordinates {
        (100,14.3)
        (1000,142.9)
        (10000,1428.6)
        (20000,2857.1)
        (chord,2857.1)
        (drums,555.6)
        (example,2.4)
      };
    \end{axis}
  \end{tikzpicture}
  \caption{Source-to-event expansion in the benchmark corpus.}\label{fig:compression-chart}
\end{figure}

This metric should not be overinterpreted.
More generated events per source line is useful when the source expresses meaningful repetition, but it can also produce
uninteresting repetition.
For a bachelor thesis, the metric is still valuable because it connects the language design to a measurable output.
The compiler removes repetitive manual event placement from the source program and makes the expansion explicit in the
MIDI artifact.

\section{Current Limits and Failure Behavior}\label{sec:limits}

The most important current limit is the event cap in the lowerer.
\texttt{MAX\_EVENTS} is 20,000.
A program that generates 20,001 note events fails during lowering with the diagnostic
\texttt{program exceeds 20000 event limit (has 20001)}.
This is a deliberate safety boundary in the implementation.
It prevents runaway programs from producing very large event streams, but it also means the thesis should not claim that
Motivo currently supports arbitrarily large scores.

The benchmark also revealed a language consistency issue in an old example file.
The current parser does not accept a key-signature header, and key signatures are not part of the current language.
The temporary benchmark copy of the example removed that unsupported line.
The thesis therefore treats key-signature and harmonic-context support as future work, not as implemented behavior.

Single-run memory checks with macOS \texttt{/usr/bin/time -l} reported maximum resident set sizes in the range of
approximately 5.2 MB to 7.8 MB for the largest synthetic cases.
This is useful as an informal local observation, but the time benchmark is the stronger measurement because it was run
with repeated trials.
The memory result should be described cautiously if used outside this chapter.

\subsection{Memory Footprint}\label{detail-subsec:memory-footprint}

The value-oriented design of the AST and IR ensures that the compiler's memory footprint remains minimal.
The use of \texttt{std::variant} supports spatial locality in the CPU cache during AST traversal.
The intermediate representation holds only the necessary atomic data (pitch, start beat, duration, velocity) packed
tightly into \texttt{struct}s. Consequently, even a composition unrolled into large event streams within the
configured implementation limits will consume only a few
megabytes of RAM, allowing the compiler to run seamlessly on low-end hardware or within memory-constrained Docker
containers.

\section{Testing Evidence}\label{sec:test-evidence}

The repository contains a broad test suite across compiler, server, and client components.
The native compiler build used during earlier local inspection exposed 582 CTest-discovered tests.
The server Vitest report in the workspace recorded 20 passing tests across 10 suites, and the client Vitest report
recorded 54 passing tests across 38 suites.
The same reports showed server line coverage of 97.95\% and client line coverage of 89.46\%.
These values should be regenerated before final submission, but they provide a useful snapshot of the current testing
discipline.

Coverage is not a proof of correctness.
It does not show whether the language is pleasant to use or whether every musical edge case has been modeled.
It does show that the implementation is treated as a software system with repeatable checks.
For this thesis, that matters because Motivo is evaluated both as a language design and as a full-stack engineering
project.

\section{Comparative Interpretation}\label{sec:comparative-interpretation}

To fully evaluate the utility and positioning of Motivo, it must be contextualized within the existing
ecosystem of music programming tools.
This section provides a comparative analysis against two distinct paradigms: live-coding environments (exemplified by
Sonic Pi) and low-level scripting libraries (such as Python's \texttt{mido}).

\subsection{Comparison with Sonic Pi}\label{detail-subsec:comparison-with-sonic-pi}

Sonic Pi is a highly successful live-coding language embedded in Ruby, designed primarily for performance and education.
While both Sonic Pi and Motivo allow music to be expressed as text, their underlying design philosophies lead
to different compositional workflows.

\subsubsection{Execution Model}
Sonic Pi utilizes a real-time, interpreted execution model.
The musician writes scripts that trigger sounds ``on the fly'' through the SuperCollider synthesis engine.
Time is managed explicitly by the programmer using the \texttt{sleep} function, which halts the execution thread.

In contrast, Motivo is fully compiled and ``zero-runtime''.
There is no concept of a thread sleeping; instead, the \texttt{play} and \texttt{rest} commands define a mathematical
relationship on a static timeline.
While Sonic Pi excels in improvisational settings where the code is meant to be modified while it runs,
Motivo is oriented toward \textit{structural composition}, where the composer needs a complex arrangement to
resolve into a consistently quantized MIDI file every time it is compiled.

\subsubsection{Structural Abstraction in Practice}
In Sonic Pi, polyphony and parallel execution are achieved by spawning new threads (e.g., using \texttt{in\_thread}).
This can lead to synchronization issues if the composer does not meticulously manage the timing of each thread.
Motivo resolves parallelism through the \texttt{voice} construct.
Because voices are unrolled at compile-time by isolating the internal cursor, the compiler supports that parallel
musical lines are aligned according to the compiler cursor model at their inception without requiring manual
thread management or timing arithmetic from the composer.

\subsection{Comparison with Raw MIDI Scripting}\label{detail-subsec:comparison-with-raw-midi-scripting}

For offline composition, many developers use general-purpose programming languages combined with MIDI libraries, such as
Python's \texttt{mido} or C++'s \texttt{RtMidi}.
While these tools offer absolute control over the output, they operate at an exceptionally low level of abstraction.

\subsubsection{Event Placement and Musical Intent}
When using a library like \texttt{mido}, the composer must manually construct \texttt{Note-On} and \texttt{Note-Off}
messages.
Crucially, because the MIDI file format relies on delta-times, the composer is forced to manually calculate the tick
difference between every sequential event across all tracks.

Motivo abstracts this entirely.
A composer writing \texttt{play (C4, E4, G4):2;} in Motivo communicates their \textit{musical intent}---playing a
C-Major triad for two beats.
Motivo's lowering engine and backend automatically handle the translation into absolute beats, the flattening into
individual notes, the sorting by absolute tick, and the final variable-length delta-time encoding.

\subsubsection{Boilerplate and Expressiveness}
A simple four-beat loop of a hi-hat playing eighth notes requires a significant amount of boilerplate in Python:
creating a file, creating a track, appending tempo metadata, and writing a loop that interleaves \texttt{Note-On} and
\texttt{Note-Off} messages with explicit delta times.

Motivo achieves the same result with minimal syntax:
\begin{lstlisting}[language=Motivo,label={lst:midi-scripting-comparison}]
tempo 120;
track using drums {
    loop (8) { play hihat :0.5; }
}
\end{lstlisting}

This contrast in verbosity illustrates Motivo's intended role as a domain-specific tool.
By embedding musical semantics (like default durations, rest tracking, and automatic track/channel assignment) directly
into the language design, the compiler reduces cognitive friction and allows the user to operate strictly within the
musical domain.

\section{Evaluation Summary}\label{sec:evaluation-summary}

The evaluation supports three bounded conclusions.
First, Motivo's compiler can expand compact source rules into much larger MIDI event streams, as shown by the synthetic
source-to-event benchmarks.
Second, the current Release build compiles the tested programs quickly on the local machine, with the largest benchmark
cases finishing around 6 ms.
Third, the implemented system includes testing and CI coverage across compiler, API, client, and paper workflows.

The evaluation also establishes limits.
The benchmark corpus is synthetic, the hardware is a single Apple M3 Pro machine, the compiler has a 20,000-note-event
safety cap, and no user study was performed.
These limits do not weaken the thesis; they make the claims more precise.
The project demonstrates a working compiler-centered music tool and measures its current behavior without claiming more
than the implementation can support.
